\bibitem{AMChen}
X.-W. Chen,
\emph{Module factorizations},
arXiv:2402.11613v3, 27 August 2025.
\url{https://arxiv.org/abs/2402.11613}.
Theorem 2.10はregular normal元によるGorenstein-projective dimensionのchange-of-rings、
Corollary 2.11は有限ambient projective dimensionが一になること、
Theorem 7.2は有限projective MFとambientで有限projective dimensionを持つ
totally reflexive加群のstable categoryとの同値を与える。
本文のgraded dg版はexact圏とdg商を通じて別途構成しており、
同論文にそのままの形で記載された定理として引用してはいない。


\bibitem{AMSasha}
Sasha,
Answer to \emph{Is the property of being a dg generator open?},
MathOverflow, 5 December 2016.
\url{https://mathoverflow.net/questions/256391/is-the-property-of-being-a-dg-generator-open}.
$\mathbb P^1$上の$\mathcal O\oplus\mathcal O(-1)[1]$が点のskyscraperへ変形して
生成性を失う例。本文では同じ例を一パラメータのstrict dg代数として書き、
全Artin切断とgenericでのsmoothnessを直接検証する。

\bibitem{AMStacks}
The Stacks Project Authors,
\emph{The Stacks Project}.
Tags \href{https://stacks.math.columbia.edu/tag/0BQT}{0BQT}
(Remark 36.16.4, projective space内の準コンパクト局所閉部分の有限線束生成子)、
\href{https://stacks.math.columbia.edu/tag/09M1}{09M1}
(Proposition 36.17.1, compactとperfectの同値)。
2026年11月10日確認。
